%% toc: P2 9 9.1,9.2,9.3,9.4
\chapter{The Syllogistic Calculi and Two Historical Reconstructions}
\label{sec:p02-w02-p02-001-historical-contrast}

The calculi \(\mathsf S\) and \(\mathsf S^\dagger\) of
Chapter~\ref{sec:p02-syllogistic-inference} are modern systems. They are
complete for a set semantics in which the universe is nonempty but a
predicate extension may be empty. This chapter states, assumption by
assumption, where those systems depart from two historical reconstructions
of Aristotle's logic: George Boger's account of Aristotle's underlying
logic, presented in Chapter~\ref{sec:p13-src063-0103} and
Section~\ref{sec:p13-boger-deduction}, and John Woods and Andrew Irvine's
account of Aristotle's early logic
\cite[pp.~27--99]{WoodsIrvine2004EarlyLogic}, which this handbook cites
directly. Neither reconstruction is identified with \(\mathsf S\) or
\(\mathsf S^\dagger\), and neither historical chapter is restated here.

\emph{Inputs.} The languages, structures and consequence relation of
Definitions~\ref{def:p02-syllogistic-language}
and~\ref{def:p02-syllogistic-semantics}; the formula complement of
Definition~\ref{def:p02-syllogistic-identifications}; the completeness
theorems~\ref{thm:p02-syllogistic-sdagger-completeness}
and~\ref{thm:p02-syllogistic-s-completeness}; and the comparison
conditions \(\alpha,\iota,\varepsilon,\omega\) of
Definition~\ref{def:p13-boger-guarded-bridge} with
Theorem~\ref{thm:p13-boger-existence-bridge}.

\emph{Conclusions.} Under an existence-committing evaluation of the four
categorical patterns, exactly the universal affirmative and the partial
privative change truth value on an empty subject
(Theorem~\ref{thm:p02-w02-p02-001-agreement}). Contrariety, subcontrariety
and subalternation hold under that evaluation and fail in \(\mathcal S\)
(Theorem~\ref{thm:p02-w02-p02-001-opposition}). The existence-committing
partial privative cannot be expressed by any set of
\(\mathcal S^\dagger\)-formulas
(Theorem~\ref{thm:p02-w02-p02-001-o-inexpressible}), but consequence under
that evaluation is decided by finitely many \(\mathsf S\)-refutations
(Theorem~\ref{thm:p02-w02-p02-001-reduction}). Finally, derivability in
\(\mathsf S\) and \(\mathsf S^\dagger\) has the structural properties that
Woods and Irvine deny to syllogisms
(Theorem~\ref{thm:p02-w02-p02-001-structural}).

\section{Differences from Boger's reconstruction}
\label{sec:p02-w02-p02-001-boger}

Boger's categorical language and its semantics are presented in
Sections~\ref{sec:p13-boger-categorical-grammar}
and~\ref{sec:p13-boger-semantics}, and his deduction apparatus in
Section~\ref{sec:p13-boger-deduction}. The comparison below names each
assumption on which \(\mathcal S\) and \(\mathcal S^\dagger\) differ.

\begin{longtable}{@{}p{0.17\linewidth}p{0.37\linewidth}p{0.37\linewidth}@{}}
\caption{Boger's reconstruction and the modern calculi}
\label{tab:p02-w02-p02-001-boger}\\
\hline
Assumption & Boger's reconstruction & \(\mathcal S\), \(\mathcal S^\dagger\) \\
\hline
\endfirsthead
\hline
Assumption & Boger's reconstruction & \(\mathcal S\), \(\mathcal S^\dagger\) \\
\hline
\endhead
Word order &
A categorical sentence combines a term, one of four logical constants and a
second term in that fixed order; the first term is the predicate and the
second the subject, as in \(AaB\)
\cite[pp.~130--131]{Boger2004UnderlyingLogic}. &
\(\forall(l,m)\) and \(\exists(l,m)\) take the subject literal first. If
atoms \(a,b\) interpret the terms \(A,B\), the pattern \(AaB\) is read as
\(\forall(b,a)\), not as \(\forall(a,b)\). \\
Identical terms &
CSFR2 requires the two terms of a categorical sentence to differ; note~29
records that Aristotle nevertheless recognises identity elsewhere
\cite[p.~131]{Boger2004UnderlyingLogic}. &
Rule T derives \(\forall(p,p)\) from no premises, and
\(\exists(p,p)\) is the standard witness premise. \\
Term negation &
The formation rules CSFR1--CSFR4 provide no operation forming the complement
of a term \cite[pp.~130--131]{Boger2004UnderlyingLogic}. In \emph{Prior
Analytics} A46, \enquote{is not-good} is classed as an affirmation, not as
the denial of \enquote{is good} \cite[p.~132]{Boger2004UnderlyingLogic};
indefinite nouns are discussed in connection with \emph{On Interpretation}
\cite[pp.~124--125]{Boger2004UnderlyingLogic}. &
Every atom has a complement literal \(\bar p\) interpreted as the set
complement. \(\mathcal S^\dagger\) allows it in subject position, and the
identification of \(\forall(l,m)\) with \(\forall(\bar m,\bar l)\) uses it. \\
Empty subjects &
Every affirmation about a nonexistent subject is false and every denial
true; with an existent subject there are no empty classes
\cite[pp.~159--163]{Boger2004UnderlyingLogic}. &
A predicate extension may be empty; \(\forall(p,q)\) is then true and
\(\exists(p,\bar q)\) false. \\
Truth conditions &
Boger declines to identify the extensional class reading with Aristotle's
own and prefers an intensional or possessional one: \(A\) belongs to every
\(B\) when every individual \(B\) has property \(A\)
\cite[pp.~164--165]{Boger2004UnderlyingLogic}. &
Purely extensional: inclusion and nonempty intersection of subsets of the
universe. \\
Rule base &
Three conversion rules, fourteen two-premise syllogism rules, and direct and
indirect (reductio) deduction \cite[p.~166]{Boger2004UnderlyingLogic}. &
The rules D1--D3, B, A, T, I and X (with N in \(\mathsf S^\dagger\)); no
premise is ever discharged. \\
\hline
\end{longtable}

The intensional reading in the fifth row has no counterpart in a structure
of Definition~\ref{def:p02-syllogistic-semantics}. The results below
therefore compare the modern calculi with the extensional comparison
semantics of Definition~\ref{def:p13-boger-guarded-bridge}, which isolates
the existence policy of the fourth row. They are statements about that
comparison semantics, not about Boger's historical interpretation as a
whole.

\begin{hbdefinition}{def:p02-w02-p02-001-patterns}{Patterns and existence-committing evaluation}
For distinct atoms \(p,q\in P\), the four \emph{patterns} with predicate
\(q\) and subject \(p\) are
\[
 q\mathbin{\mathsf a}p,\qquad q\mathbin{\mathsf e}p,\qquad
 q\mathbin{\mathsf i}p,\qquad q\mathbin{\mathsf o}p,
\]
written predicate first as in Boger's \(AaB\). Patterns with identical
terms are not formed. Let \(\mathfrak A\) be a structure of
Definition~\ref{def:p02-syllogistic-semantics}, and put
\(S=p^{\mathfrak A}\), \(P=q^{\mathfrak A}\). Define
\[
 \mathfrak A\models_{\mathrm B}q\mathbin{\mathsf a}p
 \iff\alpha(S,P),\qquad
 \mathfrak A\models_{\mathrm B}q\mathbin{\mathsf e}p
 \iff\varepsilon(S,P),
\]
\[
 \mathfrak A\models_{\mathrm B}q\mathbin{\mathsf i}p
 \iff\iota(S,P),\qquad
 \mathfrak A\models_{\mathrm B}q\mathbin{\mathsf o}p
 \iff\omega(S,P),
\]
with \(\alpha,\iota,\varepsilon,\omega\) as in
Definition~\ref{def:p13-boger-guarded-bridge}. For a set \(\Gamma\) of
patterns and a pattern \(\varphi\), \(\Gamma\models_{\mathrm B}\varphi\)
means that every structure satisfying every member of \(\Gamma\) under
\(\models_{\mathrm B}\) satisfies \(\varphi\). The \emph{contradictory}
\(\varphi^{\mathrm c}\) of a pattern keeps both terms in place and exchanges
\(\mathsf a\) with \(\mathsf o\) and \(\mathsf e\) with \(\mathsf i\).
The \emph{\(\mathcal S\)-reading} \(\sigma\) is
\[
 \sigma(q\mathbin{\mathsf a}p)=\forall(p,q),\quad
 \sigma(q\mathbin{\mathsf e}p)=\forall(p,\bar q),\quad
 \sigma(q\mathbin{\mathsf i}p)=\exists(p,q),\quad
 \sigma(q\mathbin{\mathsf o}p)=\exists(p,\bar q).
\]
\end{hbdefinition}

Each \(\sigma(\varphi)\) is an \(\mathcal S\)-formula, and
\(\sigma(\varphi^{\mathrm c})=\overline{\sigma(\varphi)}\) by
Definition~\ref{def:p02-syllogistic-identifications}. The universe of
\(\mathfrak A\) is nonempty in both evaluations, as
Definition~\ref{def:p13-boger-guarded-bridge} also requires.

\begin{hbtheorem}{thm:p02-w02-p02-001-agreement}{proved=annex:p02-w02-p02-001-agreement}{Where the two evaluations agree}
Let \(\varphi\) be a pattern with subject \(p\), and let \(\mathfrak A\)
be a structure.
\begin{enumerate}
\item If \(p^{\mathfrak A}\ne\varnothing\), then
 \(\mathfrak A\models_{\mathrm B}\varphi\) if and only if
 \(\mathfrak A\models\sigma(\varphi)\).
\item If \(p^{\mathfrak A}=\varnothing\), then
 \(\mathfrak A\models_{\mathrm B}\varphi\) exactly when \(\varphi\) is an
 \(\mathsf e\)- or \(\mathsf o\)-pattern, whereas
 \(\mathfrak A\models\sigma(\varphi)\) exactly when \(\varphi\) is an
 \(\mathsf a\)- or \(\mathsf e\)-pattern.
\end{enumerate}
Hence \(\mathsf e\)- and \(\mathsf i\)-patterns agree with their
\(\mathcal S\)-readings in every structure, and \(\mathsf a\)- and
\(\mathsf o\)-patterns disagree with them exactly when the subject is empty.
\end{hbtheorem}

\begin{hbtheorem}{thm:p02-w02-p02-001-opposition}{proved=annex:p02-w02-p02-001-opposition}{Opposition under both policies}
Let \(p,q\) be distinct atoms.
\begin{enumerate}
\item Under both evaluations each pattern and its contradictory have
 exactly one true member: \(\mathfrak A\models_{\mathrm B}\varphi\) if and
 only if \(\mathfrak A\not\models_{\mathrm B}\varphi^{\mathrm c}\), and
 \(\mathfrak A\models\sigma(\varphi)\) if and only if
 \(\mathfrak A\not\models\sigma(\varphi^{\mathrm c})\).
\item No structure satisfies both \(q\mathbin{\mathsf a}p\) and
 \(q\mathbin{\mathsf e}p\) under \(\models_{\mathrm B}\). A structure
 satisfies both \(\forall(p,q)\) and \(\forall(p,\bar q)\) if and only if
 \(p^{\mathfrak A}=\varnothing\).
\item Every structure satisfies \(q\mathbin{\mathsf i}p\) or
 \(q\mathbin{\mathsf o}p\) under \(\models_{\mathrm B}\). A structure
 satisfies neither \(\exists(p,q)\) nor \(\exists(p,\bar q)\) if and only if
 \(p^{\mathfrak A}=\varnothing\).
\item \(q\mathbin{\mathsf a}p\models_{\mathrm B}q\mathbin{\mathsf i}p\) and
 \(q\mathbin{\mathsf e}p\models_{\mathrm B}q\mathbin{\mathsf o}p\), but
 \(\forall(p,q)\not\models\exists(p,q)\) and
 \(\forall(p,\bar q)\not\models\exists(p,\bar q)\).
\end{enumerate}
\end{hbtheorem}

The conversions that Boger's deduction system licenses, and their dependence
on an existent subject, are proved in
Lemma~\ref{thm:p13-boger-conversions}. In \(\mathcal S\), the
counterpart of conversion per accidens fails for the same reason as
subalternation: an empty \(p\) makes \(\forall(p,q)\) true and
\(\exists(q,p)\) false.

\begin{hbtheorem}{thm:p02-w02-p02-001-o-inexpressible}{proved=annex:p02-w02-p02-001-o-inexpressible}{The partial privative is inexpressible}
For distinct atoms \(p,q\) and every structure \(\mathfrak A\),
\[
 \mathfrak A\models_{\mathrm B}q\mathbin{\mathsf a}p\iff
 \mathfrak A\models\{\forall(p,q),\exists(p,p)\},\qquad
 \mathfrak A\models_{\mathrm B}q\mathbin{\mathsf e}p\iff
 \mathfrak A\models\forall(p,\bar q),
\]
\[
 \mathfrak A\models_{\mathrm B}q\mathbin{\mathsf i}p\iff
 \mathfrak A\models\exists(p,q),
\]
\[
 \mathfrak A\models_{\mathrm B}q\mathbin{\mathsf o}p\iff
 \mathfrak A\models\forall(p,\bar p)\quad\text{or}\quad
 \mathfrak A\models\exists(p,\bar q).
\]
There is, however, no set \(\Theta\subseteq\mathcal S^\dagger\) such that
for every structure \(\mathfrak A\),
\(\mathfrak A\models\Theta\) if and only if
\(\mathfrak A\models_{\mathrm B}q\mathbin{\mathsf o}p\).
\end{hbtheorem}

The partial privative therefore enters a translation into the modern
calculi only as a two-way case distinction. The next result shows that
this case distinction suffices.

\begin{hbdefinition}{def:p02-w02-p02-001-selection}{Translations under a selection}
Let \(\Delta\) be a finite set of patterns. A \emph{selection} \(s\) for
\(\Delta\) assigns to each \(\mathsf o\)-pattern \(q\mathbin{\mathsf o}p\in\Delta\)
one of the two formulas \(\forall(p,\bar p)\), \(\exists(p,\bar q)\). Put
\[
 \tau_s(q\mathbin{\mathsf a}p)=\{\forall(p,q),\exists(p,p)\},\quad
 \tau_s(q\mathbin{\mathsf e}p)=\{\forall(p,\bar q)\},\quad
 \tau_s(q\mathbin{\mathsf i}p)=\{\exists(p,q)\},
\]
\(\tau_s(q\mathbin{\mathsf o}p)=\{s(q\mathbin{\mathsf o}p)\}\), and
\(\tau_s(\Delta)=\bigcup_{\psi\in\Delta}\tau_s(\psi)\). Every member of
\(\tau_s(\Delta)\) is an \(\mathcal S\)-formula. When \(\Delta\) contains \(k\) distinct
\(\mathsf o\)-patterns there are exactly \(2^k\) selections.
\end{hbdefinition}

\begin{hbtheorem}{thm:p02-w02-p02-001-reduction}{proved=annex:p02-w02-p02-001-reduction}{A reduction to \(\mathsf S\)}
Let \(\Gamma\) be a finite set of patterns, \(\varphi\) a pattern, and
\(\Delta=\Gamma\cup\{\varphi^{\mathrm c}\}\). Fix an atom \(r\). Then
\(\Gamma\models_{\mathrm B}\varphi\) if and only if, for every selection
\(s\) for \(\Delta\),
\[
 \tau_s(\Delta)\vdash_{\mathsf S}\exists(r,\bar r).
\]
\end{hbtheorem}

\begin{hbmethod}{def:p02-w02-p02-001-procedure}{A two-policy inference test}
Given premises \(\Gamma\) and a proposed conclusion \(\varphi\), all
patterns with distinct terms:
\begin{enumerate}
\item \emph{Modern policy.} Test
 \(\sigma[\Gamma]\models\sigma(\varphi)\) with the saturation method of
 Definition~\ref{def:p02-syllogistic-saturation}. When it fails, retain an
 empty extension in the countermodel search.
\item \emph{Existence-committing policy.} Form
 \(\Delta=\Gamma\cup\{\varphi^{\mathrm c}\}\). For each of the \(2^k\)
 selections, derive an absurdity from \(\tau_s(\Delta)\) in \(\mathsf S\),
 or exhibit a structure satisfying \(\tau_s(\Delta)\). By
 Theorem~\ref{thm:p02-w02-p02-001-reduction}, the inference holds exactly
 when every selection is refuted; a structure satisfying one
 \(\tau_s(\Delta)\) is a countermodel under \(\models_{\mathrm B}\).
\item \emph{Compare.} By Theorem~\ref{thm:p02-w02-p02-001-agreement}, a
 difference between the two verdicts requires an empty subject term in the
 countermodel of the inference that failed.
\end{enumerate}
\emph{Side conditions.} The terms of every pattern are distinct atoms.
The universe is nonempty in both policies. Some is read indeterminately:
at least one, possibly all.
\end{hbmethod}

\begin{hbexample}{def:p02-w02-p02-001-darapti}{Valid only with existential import}
Let \(w,m,s\) mean whale, mammal and swimmer, and take premises
\(m\mathbin{\mathsf a}w\) and \(s\mathbin{\mathsf a}w\), with proposed
conclusion \(m\mathbin{\mathsf i}s\): some swimmer is a mammal.
There is no \(\mathsf o\)-pattern in \(\Delta\), whose contradictory member
is \(m\mathbin{\mathsf e}s\), so the only selection gives
\[
 \tau_s(\Delta)=\{\forall(w,m),\ \exists(w,w),\ \forall(w,s),\ \forall(s,\bar m)\}.
\]
D1 applied to \(\forall(w,s)\) and \(\exists(w,w)\) gives
\(\exists(w,s)\). D2 applied to \(\exists(w,s)\) and \(\forall(w,m)\) gives
\(\exists(m,s)\), identified with \(\exists(s,m)\). Its formula complement
is \(\forall(s,\bar m)\), so X yields \(\exists(w,\bar w)\). The inference
is valid under \(\models_{\mathrm B}\).

Under the modern policy the premises are \(\forall(w,m)\) and
\(\forall(w,s)\). The structure with universe \(\{c\}\) and
\(w^{\mathfrak A}=m^{\mathfrak A}=s^{\mathfrak A}=\varnothing\) satisfies
both and falsifies \(\exists(s,m)\). The verdicts differ exactly because
of the empty subject \(w\), as step~3 of the method predicts. Supplying the
premise \(\exists(w,w)\) restores the conclusion in \(\mathsf S\) by the
same two steps.
\end{hbexample}

\begin{hbexample}{def:p02-w02-p02-001-baroco}{A case split for a partial privative premise}
Take premises \(m\mathbin{\mathsf a}s\) (every swimmer is a mammal) and
\(m\mathbin{\mathsf o}w\), with proposed conclusion
\(s\mathbin{\mathsf o}w\). Then \(\Delta\) adds \(s\mathbin{\mathsf a}w\),
and the two selections give
\[
 \{\forall(s,m),\exists(s,s),\forall(w,\bar w),\forall(w,s),\exists(w,w)\}
\]
and
\[
 \{\forall(s,m),\exists(s,s),\exists(w,\bar m),\forall(w,s),\exists(w,w)\}.
\]
In the first set \(\forall(w,\bar w)\) and \(\exists(w,w)\) are formula
complements, and X yields an absurdity. In the second, B applied to
\(\forall(w,s)\) and \(\forall(s,m)\) gives \(\forall(w,m)\), whose
complement \(\exists(w,\bar m)\) is present, and X again yields an
absurdity. Both selections are refuted, so the inference is valid under
\(\models_{\mathrm B}\). The first selection is the empty-subject case, in
which the existence-committing partial privative premise is true without a
witness.
\end{hbexample}

\begin{hbwarning}{def:p02-w02-p02-001-limits}{What the comparison does not capture}
The evaluation \(\models_{\mathrm B}\) is extensional. It does not
represent Boger's intensional or possessional truth conditions, essential
versus accidental attribution, or the determinate reading of some, under
which a partial sentence also requires a member lacking the property
\cite[pp.~147--148, 164--165]{Boger2004UnderlyingLogic}. A validity verdict
under \(\models_{\mathrm B}\) is therefore not a verdict about Boger's
reconstruction in every respect. Theorem~\ref{thm:p02-w02-p02-001-o-inexpressible}
also limits the reduction: the number of \(\mathsf S\)-refutations doubles
with each distinct partial privative in \(\Delta\).
\end{hbwarning}

\section{Differences from Woods and Irvine's reconstruction}
\label{sec:p02-w02-p02-001-woods-irvine}

Woods and Irvine reconstruct the theory of syllogisms in the broad sense
that, on their reading, the \emph{Topics} and \emph{On Sophistical
Refutations} contain and the \emph{Prior Analytics} presupposes
\cite[pp.~38, 51--56]{WoodsIrvine2004EarlyLogic}. Their assumptions differ
from those of \(\mathsf S\) and \(\mathsf S^\dagger\) as follows.

\begin{itemize}
\item \emph{Two relations, not one.} They separate validity, on which the
 premises necessitate the conclusion, from the property of being a
 syllogism, which restricts validity. They attribute to validity the
 structural properties of reflexivity, transitivity and monotonicity, and
 adopt a classical theory of validity as an interpretive policy; they deny
 reflexivity and monotonicity of syllogisms
 \cite[pp.~49--51, 59]{WoodsIrvine2004EarlyLogic}.
 \(\vdash_{\mathsf S}\) and \(\vdash_{\mathsf S^\dagger}\) are complete
 for a single consequence relation and contain no restriction of this kind.
\item \emph{Conditions on syllogisms.} A syllogism is a valid argument all
 of whose statements are propositions in Aristotle's technical sense, and
 it satisfies Min (no premise beyond those needed for validity), Non-Circ
 (the conclusion repeats no premise) and, as a provisional condition, Prem+
 (more than one premise) \cite[pp.~40, 54--56]{WoodsIrvine2004EarlyLogic}.
\item \emph{Inconsistent premises.} From Non-Circ and the principle of
 argumental conversion they conclude that no argument with expressly
 inconsistent premises is a syllogism, and they describe the resulting
 logic as paraconsistent \cite[pp.~64--65]{WoodsIrvine2004EarlyLogic}.
 Rule X of both sheets derives every formula of the fragment from
 \(\psi\) and \(\bar\psi\).
\item \emph{Propositions.} A proposition says one thing of one thing and,
 apart from negation, contains no connective; their example
 \enquote{No men are non-animals} counts as a proposition, so negation
 inside the predicate is permitted
 \cite[pp.~39--40, 65]{WoodsIrvine2004EarlyLogic}. \(\mathcal S^\dagger\)
 also allows complemented literals, in either position, while excluding
 conjunction and disjunction (Definition~\ref{def:p02-syllogistic-language}).
\item \emph{Direct and hypothetical syllogisms.} In a direct syllogism every
 line other than the conclusion is a premise; a hypothetical
 (\emph{per impossibile}) syllogism introduces a hypothesis that is later
 cancelled \cite[p.~52]{WoodsIrvine2004EarlyLogic}. \(\mathsf S\) and
 \(\mathsf S^\dagger\) derivations are direct in this sense.
\item \emph{An existential example.} Among valid arguments, their
 Table~2, row~2 lists \enquote{All men are mortal; All ducks are quackers;
 therefore, some men are mortal}
 \cite[p.~63]{WoodsIrvine2004EarlyLogic}. Under the modern policy the
 corresponding \(\forall(p,q)\models\exists(p,q)\) fails
 (Theorem~\ref{thm:p02-w02-p02-001-opposition}, item~4), so the example
 is valid only with an existence assumption on the subject term.
\end{itemize}

\begin{hbtheorem}{thm:p02-w02-p02-001-structural}{proved=annex:p02-w02-p02-001-structural}{Structural properties of the modern derivability relations}
Let \(\vdash\) be \(\vdash_{\mathsf S}\) on \(\mathcal S\) or
\(\vdash_{\mathsf S^\dagger}\) on \(\mathcal S^\dagger\), and let all
formulas belong to the corresponding fragment.
\begin{enumerate}
\item \emph{Reflexivity.} If \(\theta\in\Theta\), then
 \(\Theta\vdash\theta\).
\item \emph{Monotonicity.} If \(\Theta\vdash\theta\) and
 \(\Theta\subseteq\Theta'\), then \(\Theta'\vdash\theta\).
\item \emph{Transitivity.} If \(\Theta\vdash\theta_i\) for
 \(i=1,\ldots,n\) and \(\Theta\cup\{\theta_1,\ldots,\theta_n\}\vdash\theta\),
 then \(\Theta\vdash\theta\).
\item \emph{Explosion.} \(\{\psi,\bar\psi\}\vdash\varphi\) for all
 \(\psi,\varphi\).
\item \emph{Premise-free derivation.} \(\varnothing\vdash\forall(p,p)\) for
 every atom \(p\).
\item \emph{No subalternation.} For distinct atoms \(p,q\),
 \(\{\forall(p,q)\}\nvdash\exists(p,q)\).
\end{enumerate}
\end{hbtheorem}

By items 1 and~2, \(\{\forall(p,q),\forall(o,\bar o)\}\vdash\forall(p,q)\),
an argument that fails Min and Non-Circ. By item~4,
\(\{\exists(p,p),\forall(p,\bar p)\}\vdash\forall(q,o)\), an argument with
expressly inconsistent premises. By item~5, \(\forall(p,p)\) is derivable
from no premises, the situation Woods and Irvine discuss for logical truths
\cite[pp.~60--62]{WoodsIrvine2004EarlyLogic}. The modern calculi
therefore model a consequence relation with the structural profile that
Woods and Irvine assign to validity, not the restricted relation they assign
to syllogisms. The completeness theorems of
Chapter~\ref{sec:p02-syllogistic-inference} concern that unrestricted
relation; they establish nothing about its restriction by Min, Non-Circ and
Prem+.
\index{syllogistic logic!existential import}
\index{syllogistic logic!historical reconstructions}
\index{existential import!Boger's reconstruction}
